% ECONOMICS_PROBLEM: {"id": "D82-56", "title": "Robust mechanism design", "jel_code": "D82", "source_id": "OP-0005"}
% !TeX program = xelatex
\documentclass[11pt,a4paper]{article}
\usepackage[margin=23mm]{geometry}
\usepackage{fontspec}
\setmainfont{Latin Modern Roman}
\usepackage{amsmath,amssymb,mathtools}
\usepackage{xcolor,enumitem,booktabs,longtable,array,xurl,hyperref}
\definecolor{muted}{HTML}{53636C}
\hypersetup{colorlinks=true,urlcolor=blue,linkcolor=blue}
\setlength{\parindent}{0pt}
\setlength{\parskip}{5pt}
\newcommand{\R}{\mathbb R}
\newcommand{\E}{\mathbb E}
\newcommand{\N}{\mathbb N}
\newcommand{\ind}{\mathbf 1}
\newcommand{\Var}{\operatorname{Var}}
\newcommand{\Cov}{\operatorname{Cov}}
\newcommand{\supp}{\operatorname{supp}}
\DeclareMathOperator*{\argmin}{arg\,min}
\DeclareMathOperator*{\argmax}{arg\,max}
\newcommand{\entrymeta}[3]{{\sffamily\small\color{muted}\textbf{\texttt{#1}}\quad|\quad #2\par #3\par}}
\newcommand{\Problemref}[1]{\texttt{#1}}
\newcommand{\JELref}[2]{\texttt{#2} (JEL \texttt{#1})}
\begin{document}
\section*{Robust mechanism design}
\textbf{Problem ID:} \texttt{D82-56}\quad \textbf{JEL:} \texttt{D82}\par
% BEGIN ECONOMICS STATEMENT
\label{problem:OP-0005}
\label{atlas:prob:M5}

\noindent\textbf{Original atlas ID:} M5. \textbf{Formulation:} editorial specification of a research family.

\subsection*{Definitions and assumptions}

Specify finite player, payoff-state, type, and action sets. Let $\mathfrak I$ be a nonempty, declared collection of common-prior information structures, each giving a joint distribution of payoff states and private signals. Let $\mathfrak M$ be an admissible class of mechanisms with finite message sets, feasible outcome lotteries, and bounded transfers. For $(m,I)$ let $\mathcal E(m,I)$ denote its Bayes--Nash equilibria. Designer payoff $W(m,I,e)$ is the expected value of a fixed bounded welfare or revenue function. Participation, allocation feasibility, and any budget constraints must hold in the precise interim or ex-post sense chosen for every admissible structure and equilibrium. Write $\mathfrak M_F$ for mechanisms satisfying these robust constraints, and assume this set is nonempty.

\subsection*{Formal research question}

Characterize and compute the robust value
\[
V^{\rm rob}=\sup_{m\in\mathfrak M_F}\inf_{I\in\mathfrak I}\inf_{e\in\mathcal E(m,I)}W(m,I,e).
\]
For every requested $\varepsilon>0$, can one construct an admissible $m_\varepsilon$ whose worst-case payoff is at least $V^{\rm rob}-\varepsilon$, with a verifiable guarantee? Investigate how the guarantee changes as $\mathfrak I$ is enlarged, and identify informational restrictions under which simpler mechanisms remain optimal or approximately optimal. If implementation is instead required for all rationalizable outcomes, replace the equilibrium correspondence explicitly; if non-common priors are allowed, define individual beliefs and the solution concept rather than reusing the common-prior Bayes formula.

\subsection*{Interpretation and scope}

A max--min value is a mathematically specified objective, not a universal prescription. An ambiguity set allowing arbitrarily unfavorable values can make the value trivial, while a tightly estimated set can restore informative performance. The inner infimum over equilibria records pessimistic selection; choosing an advantageous equilibrium would represent a different problem. Finite message spaces do not justify assuming a revelation principle for every robustness or implementation requirement. Empirical comparisons should assess informational misspecification and realized participation, and should distinguish robustness of predictions from robustness of mechanism performance. The relevant target may instead be regret relative to a structure-specific oracle, which requires a different objective and reference mechanism class.

\subsection*{Source audit and status}

Bergemann--Morris (2005) and Carroll (2015) are verified. Carroll's result concerns uncertainty about available actions in a moral-hazard problem, so it is not a general adverse-selection solution. The shorthand ``Bergemann--Morris (2016)'' is ambiguous: both the Theoretical Economics article [S2] and the Papers and Proceedings article [S4] are verified and are listed explicitly. Bayes correlated equilibrium provides obedience-based outcome restrictions; it does not itself optimize every robust mechanism. Original link [11] correctly points to the 2005 paper. The max--min formulation above is editorial, and broad informational robustness should not be called an unresolved theorem without a narrower primitive class.

\subsection*{References}

\begin{enumerate}

\item[S1] Dirk Bergemann and Stephen Morris (2005). \emph{Robust Mechanism Design}. Econometrica 73(6), 1771--1813. \url{https://onlinelibrary.wiley.com/doi/10.1111/j.1468-0262.2005.00638.x}. \textit{Locator:} abstract and informational-robustness formulation. \textit{Establishes:} Robustness to higher-order belief specifications in mechanism design. \textit{Does not establish:} optimal design under every ambiguity set and behavioral model.

\item[S2] Dirk Bergemann and Stephen Morris (2016). \emph{Bayes Correlated Equilibrium and the Comparison of Information Structures in Games}. Theoretical Economics 11(2), 487--522. \url{https://www.econtheory.org/ojs/index.php/te/article/viewArticle/20160487}. \textit{Locator:} abstract; obedience characterization. \textit{Establishes:} An outcome characterization and comparison of information structures using Bayes correlated equilibrium. \textit{Does not establish:} a general robust mechanism optimizer.

\item[S3] Gabriel Carroll (2015). \emph{Robustness and Linear Contracts}. American Economic Review 105(2), 536--563. \url{https://www.aeaweb.org/articles?id=10.1257/aer.20131159}. \textit{Locator:} abstract; robust moral-hazard model. \textit{Establishes:} Optimal linear contracts for the specified uncertainty about available actions. \textit{Does not establish:} that linear contracts solve every robust adverse-selection problem.

\item[S4] Dirk Bergemann and Stephen Morris (2016). \emph{Information Design, Bayesian Persuasion, and Bayes Correlated Equilibrium}. American Economic Review: Papers and Proceedings 106(5), 586--591. \url{https://pubs.aeaweb.org/doi/10.1257/aer.p20161046}. \textit{Locator:} abstract. \textit{Establishes:} Information-design interpretation of Bayes correlated equilibrium. \textit{Does not establish:} a resolution of endogenous costly learning in general.

\end{enumerate}

\subsection*{Common conventions: Source mathematical conventions}

Unless an entry states otherwise, agent and firm sets are finite. State and action spaces are measurable, strategies and outcome maps are measurable, probability laws are specified, and expectations exist for the targets under discussion. Infinite-horizon discount factors lie in $(0,1)$ and discounted objectives are finite. Norms, tolerances and complexity input sizes must be specified for computational comparisons. Constants or primitives that appear locally are inputs, not empirical facts asserted by this catalogue.

\textbf{Identification.} A statistical model $m\in\mathcal M$ implies an observable law $P_m$ and target $\theta(m)$. At law $P$, its identified set is
\[
\Theta(P;\mathcal M)=\{\theta(m):m\in\mathcal M,\ P_m=P\}.
\]
Point identification holds when this set is a singleton, or equivalently when $P_{m_1}=P_{m_2}$ implies $\theta(m_1)=\theta(m_2)$ for all compatible models. Sharp bounds mean bounds attained, or approached when the boundary is not attained, by compatible models. Writing this set is a definition, not a solution: the substantive task is to establish informative restrictions and derive or estimate the set. An unrestricted-confounding model may give only trivial bounds.

\textbf{Inference.} Identification, estimability and valid uncertainty quantification are separate questions. Fix $\alpha\in(0,1)$. A confidence procedure $C_n$ over a declared law class $\mathcal P$ has asymptotic uniform coverage at level $1-\alpha$ when
\[
\liminf_{n\to\infty}\inf_{P\in\mathcal P}\Pr_P\{\theta(P)\in C_n\}\geq1-\alpha.
\]
For set-identified targets the coverage event must be defined separately, for example coverage of the full identified set. If an entry uses identification-indexed models instead of uniquely indexed laws, the same coverage convention is applied over those models. Pointwise limits do not establish uniform validity, and asymptotic validity does not guarantee finite-sample precision. Sample sizes, dependence structures and weak-identification sequences are part of each application.

\textbf{Counterfactuals.} $Y_i(a)$ is a potential outcome under a specified intervention, including its timing, implementation and versions. It is not $Y_i$ conditional on voluntarily choosing $a$. A full intervention vector permits interference; shorthand scalar treatment notation does not silently impose no interference. Assignment, selection, attrition, anticipation and measurement must be specified. A claim about an instrument's local effect is not automatically a population policy effect.

\textbf{Economic equilibrium.} A counterfactual model includes best responses, constraints, feasibility, market clearing, state transitions and expectations consistent with the stated information. When several equilibria exist, either a declared selector is an input or the target is set-valued. Comparisons across policy regimes must use the stated selection convention. Reduced-form relationships are not presumed invariant after an equilibrium-changing intervention.

\textbf{Optimization and welfare.} Optimization is over the stated feasible set. If attainment is not established, $\sup$ or $\inf$ and $\varepsilon$-optimal choices replace an asserted maximizer or minimizer. For example, with value $V=\sup_{a\in\mathcal A}W(a)<\infty$, an $\varepsilon$-optimal set is $\{a\in\mathcal A:W(a)\geq V-\varepsilon\}$ for $\varepsilon>0$. Benefits, costs, risk and health or environmental outcomes can be aggregated only using declared units and conversion weights. Interpersonal utility weights, the treatment of fiscal transfers and foreign profits, discounting, and the behavioral welfare criterion are explicit inputs. A comparative-static derivative additionally requires existence and appropriate differentiability of the selected counterfactual.

\textbf{Sources.} Local labels [S1], [S2], and so forth restart in every question. Locators identify the relevant abstract, model, section or result at the level actually checked; a bibliographic or abstract-level verification is not represented as inspection of an entire proof. Working papers are distinguished from later publications and changing versions. Source summaries are paraphrased. All URL checks and editorial formulations refer to the audit date stated above.
% END ECONOMICS STATEMENT
\end{document}
